\section{Protocol and Approximation Model}

\subsection{Authenticated Packet Layer}

Ghostext packet format is
\begin{equation}
\texttt{salt} \parallel \texttt{header\_nonce} \parallel \texttt{sealed\_header} \parallel \texttt{body\_nonce} \parallel \texttt{body\_ciphertext}.
\end{equation}
The internal header stores protocol version, crypto IDs, body length, and configuration fingerprint. Current implementation uses scrypt for key derivation, HKDF-SHA256 for key separation, and ChaCha20-Poly1305 for authenticated encryption.

\subsection{Candidate Policy}

At step $t$, raw logits define base distribution $P_t$ over vocabulary $\mathcal{V}$. Ghostext constructs candidate set $C_t$ via deterministic top-$p$ and max-$k$ truncation, then renormalizes to $\widetilde{P}_t$. If $|C_t|<2$ or step entropy is below threshold, no payload bit is consumed and the top token is emitted.

Retokenization-stability filtering then removes candidates that fail render/re-tokenize path consistency.

\subsection{Integer Quantization}

Let $N_t=|C_t|$ and total frequency budget be $F_{\mathrm{tot}}$ (default $65536$ in current implementation). For each candidate $v\in C_t$:
\begin{equation}
 f_t(v)=1+\left\lfloor \widetilde{P}_t(v)(F_{\mathrm{tot}}-N_t)\right\rfloor,
\end{equation}
with leftover units distributed by fractional remainder and token-id tie-break.

The induced quantized distribution is
\begin{equation}
Q_t(v)=\frac{f_t(v)}{F_{\mathrm{tot}}}.
\end{equation}

\subsection{Interval Update}

Given interval $[L_t,H_t)$ and cumulative boundaries from $Q_t$, selecting token $v$ updates
\begin{align}
L_{t+1}&=L_t+\left\lfloor(H_t-L_t)\,\mathrm{cdf}^-_t(v)\right\rfloor,\\
H_{t+1}&=L_t+\left\lfloor(H_t-L_t)\,\mathrm{cdf}^+_t(v)\right\rfloor.
\end{align}
A segment resolves when $H_t-L_t=1$.

\subsection{Bound for Quantization Error}

For each candidate, deterministic rounding implies
\begin{equation}
\left|Q_t(v)-\widetilde{P}_t(v)\right|\le\frac{1}{F_{\mathrm{tot}}}.
\end{equation}
Hence
\begin{equation}
\|Q_t-\widetilde{P}_t\|_1\le\frac{N_t}{F_{\mathrm{tot}}}, \quad
\mathrm{TV}(Q_t,\widetilde{P}_t)\le\frac{N_t}{2F_{\mathrm{tot}}}.
\end{equation}

\subsection{Bound for Truncation + Quantization + Stability Filtering}

Define truncated tail mass
\begin{equation}
\alpha_t = 1-\sum_{v\in C_t}P_t(v),
\end{equation}
and filtered-out mass under $Q_t$
\begin{equation}
\beta_t = \sum_{v\in C_t\setminus S_t}Q_t(v),
\end{equation}
where $S_t$ is the stability-surviving subset.

Let $R_t$ be $Q_t$ restricted to $S_t$ before renormalization. Then
$\|R_t-Q_t\|_1=\beta_t$ and, because $Q'_t=R_t/(1-\beta_t)$,
$\|Q'_t-R_t\|_1=\beta_t$. Therefore
\begin{equation}
\mathrm{TV}(Q'_t,Q_t)\le\beta_t,
\end{equation}
so the normalization effect after filtering is already covered by the same $\beta_t$ term.

Combining truncation, quantization, and filtering gives
\begin{equation}
\mathrm{TV}(Q'_t,P_t)\le\alpha_t + \frac{N_t}{2F_{\mathrm{tot}}} + \beta_t.
\label{eq:tv_total_bound}
\end{equation}
This bound is implementation-aligned and directly measurable from runtime statistics.

\subsection{Fingerprint Collision Discussion}

Configuration fingerprint uses 64-bit truncation for synchronization mismatch detection. For $n$ independent sessions,
\begin{equation}
\Pr[\mathrm{collision}]\approx\frac{n(n-1)}{2^{65}}.
\end{equation}
For $n=10^6$, this is approximately $2.7\times10^{-8}$. Packet authenticity still depends on AEAD tags; fingerprint collisions do not replace cryptographic integrity checks.

In adversarial settings, a party might search configuration variants that collide in 64 bits. This does not produce valid plaintext release without passphrase-bound AEAD success, but it can degrade diagnostic clarity by making a mismatch appear superficially aligned at the fingerprint layer. For this reason, fingerprint match should be interpreted as a fast guard, while authenticity remains cryptographic.

\subsection{Concrete Encode/Decode Examples}

\begin{table*}[t]
\caption{Two concrete end-to-end examples from current toy-backend implementation.}
\label{tab:examples}
\begin{tabular}{p{0.17\linewidth}p{0.20\linewidth}p{0.40\linewidth}p{0.18\linewidth}}
\toprule
Prompt type & Secret message & Cover text (prefix) & Outcome / token stats \\
\midrule
English walk prompt & \texttt{Meet at 7 PM by the old bridge.} & \texttt{The afternoon feels calm ... people talk softly along the street ...} & decode success; total 455; packet 428; tail 27; bits/token 2.092 \\
Chinese walk prompt (romanized message) & \texttt{Jin wan qi dian zai lao di fang jian.} & \texttt{Jin tian de feng hen qing, jie jiao de kafei dian hai liang zhe deng ...} & decode success; total 424; packet 415; tail 9; bits/token 2.226 \\
\bottomrule
\end{tabular}
\end{table*}
